% W4i44_Observables.tex
% DFD 17-Agent Research Programme --- Iteration 44 (i44)
% Agents 8 (Laboratory), 9 (Particle Physics), 10 (Astrophysics), 11 (CMB), 12 (LSS), 13 (Dark Sector), 14 (SymBoltz/Numerics)
% Theme: SymBoltz Phase 0 execution, Euclid pre-registration, Th-229 electroplating, DESI DR2 published, A_L joint
% Date: 2026-03-23

\documentclass[11pt,a4paper]{article}
\usepackage[utf8]{inputenc}
\usepackage{amsmath,amssymb,amsfonts,amsthm}
\usepackage{hyperref}
\usepackage{booktabs}
\usepackage{longtable}
\usepackage{geometry}
\usepackage{xcolor}
\usepackage{tcolorbox}
\geometry{margin=2.5cm}

\newtheorem{theorem}{Theorem}
\newtheorem{proposition}{Proposition}
\newtheorem{lemma}{Lemma}
\newtheorem{corollary}{Corollary}
\newtheorem{definition}{Definition}
\newtheorem{remark}{Remark}

\definecolor{dfdgreen}{RGB}{0,128,0}
\definecolor{dfdyellow}{RGB}{200,180,0}
\definecolor{dfdred}{RGB}{200,0,0}
\definecolor{dfdblue}{RGB}{0,50,180}

\newcommand{\GREEN}{\textcolor{dfdgreen}{\textbf{GREEN}}}
\newcommand{\YELLOW}{\textcolor{dfdyellow}{\textbf{YELLOW}}}
\newcommand{\RED}{\textcolor{dfdred}{\textbf{RED}}}
\newcommand{\NEW}{\textcolor{dfdblue}{\textbf{NEW}}}
\newcommand{\RESOLVED}{\textcolor{dfdgreen}{\textbf{RESOLVED}}}
\newcommand{\Tr}{\operatorname{Tr}}
\newcommand{\CP}{\mathbb{C}P}

\title{DFD i44: Observables --- Laboratory, Particle, Astro, CMB, LSS, Dark Sector, Numerics\\[6pt]
\large Agents 8, 9, 10, 11, 12, 13, 14\\[6pt]
\normalsize Iteration 13/20 of Quantum Gravity Cycle (i32--i51)}
\author{DFD 17-Agent Research Programme}
\date{2026-03-23}

\begin{document}
\maketitle
\tableofcontents
\newpage

%=============================================================================
\part{Agent 8: Laboratory Experiments}
\label{part:agent8}
%=============================================================================

\section{Mandate}

Update on Th-229 nuclear clock, SRF cavities, SQMS, and precision measurements. Special focus on $\dot{\alpha}/\alpha$ timeline.

\section{\NEW: Thorium-229 Electroplating Breakthrough}

\begin{tcolorbox}[colback=blue!5!white, colframe=blue!60!black,
  title={Electroplated Th-229 on Steel (ScienceDaily, January 2026)}]
A completely new method of probing the Th-229 nuclear isomeric transition has been demonstrated: by electroplating thorium onto steel substrates. This bypasses the need for delicate single-crystal CaF$_2$ hosts entirely.

\textbf{Significance}: The crystal growth bottleneck for Th-229 clocks may be removed. If electroplated samples achieve comparable frequency stability, the timeline to operational nuclear clocks could be significantly accelerated.

\textbf{Caveat}: Frequency reproducibility has not yet been demonstrated with electroplated samples. The CaF$_2$ results (220~Hz reproducibility, $1.1 \times 10^{-13}$ fractional) remain the benchmark.
\end{tcolorbox}

\section{Th-229:CaF$_2$ Frequency Reproducibility (Continued)}

The Nature (January 2026) result stands:
\begin{itemize}
\item At 195~K: reproducibility 220~Hz ($1.1 \times 10^{-13}$) over 7 months between differently doped crystals.
\item Optimal working temperature: $196 \pm 5$~K (first-order thermal sensitivity vanishes).
\item At this temperature, in-situ co-sensing reduces temperature-induced systematics below $10^{-18}$.
\end{itemize}

\section{\NEW: Fine-Structure Constant Sensitivity Coefficient}

Nature Communications (2025): Ma Yugang's team at Fudan computed $K_\alpha$ for $^{229}$Th$^{3+}$ using multiconfiguration Dirac--Hartree--Fock. The absolute sensitivity coefficient determines how precisely the Th-229 clock can constrain $\dot{\alpha}/\alpha$.

\textbf{DFD impact}: With $K_\alpha$ now known, the projected sensitivity to $\dot{\alpha}/\alpha \sim 10^{-19}$/yr by $\sim$2030 is confirmed as physically achievable, not merely an extrapolation.

\section{Revised Th-229 Timeline}

\begin{center}
\begin{tabular}{lll}
\toprule
\textbf{Year} & \textbf{Milestone} & \textbf{Status} \\
\midrule
2025 & Frequency reproducibility demonstrated & DONE \\
2026 & Electroplated samples tested & IN PROGRESS \\
$\sim$2028 & Operational solid-state Th-229 clock & ON TRACK \\
$\sim$2029 & $10^{-18}$ fractional sensitivity & ON TRACK \\
$\sim$2030--2031 & DFD-relevant $10^{-19}$/yr for $\dot{\alpha}/\alpha$ & ON TRACK \\
\bottomrule
\end{tabular}
\end{center}

$\sim$12 research teams worldwide (China, Europe, Japan, US) are independently pursuing nuclear clock development.

\section{SRF and Dark Photon Searches}

\begin{itemize}
\item \textbf{Dark SRF refined bounds} (PRL, March 2026): Improved theoretical modelling of frequency instability leads to dark-photon exclusion bounds an order of magnitude stronger than previous pathfinder results.
\item \textbf{SQMS 2.0}: \$125M renewal continues. Tunable SRF cavities under development for broader mass range in dark photon/axion searches.
\end{itemize}

\section{Agent 8 Assessment: Grade A-}

The electroplating breakthrough is a genuine advance. The $K_\alpha$ determination confirms the DFD $\dot{\alpha}/\alpha$ test is feasible. Dark SRF bounds tightened.

\section{New Literature (Agent 8)}

\begin{enumerate}
\item \textbf{Electroplated Th-229 on steel} (ScienceDaily, January 2026) --- New production method. \NEW.
\item \textbf{Nature Comms (2025)} --- $K_\alpha$ sensitivity coefficient for Th-229. \NEW.
\item \textbf{Dark SRF refined bounds} (PRL, March 2026) --- Order-of-magnitude improvement. \NEW.
\item \textbf{Nature (January 2026)} --- Frequency reproducibility at $1.1 \times 10^{-13}$.
\item \textbf{PRL 134 (March 2025)} --- Temperature sensitivity of Th-229 transition.
\item \textbf{Hudson et al.\ (December 2025)} --- IC electron M\"ossbauer spectroscopy.
\end{enumerate}


%=============================================================================
\part{Agent 9: Particle Physics}
\label{part:agent9}
%=============================================================================

\section{Mandate}

LHC Run 3 and HL-LHC timeline. Higgs mass precision update. BSM search status.

\section{HL-LHC Timeline Update}

\begin{tcolorbox}[colback=blue!5!white, colframe=blue!60!black,
  title={HL-LHC Schedule (March 2026)}]
\begin{itemize}
\item \textbf{LHC Run 3}: Extended until \textbf{July 2026}.
\item \textbf{Long Shutdown 3 (LS3)}: Begins July 2026.
\item \textbf{High-intensity beam test}: End of 2026 run (HL-LHC-type beams in LHC).
\item \textbf{MQXFB quadrupole production}: $\sim$2/3 of coils completed; all components to CERN by end 2026.
\item \textbf{HL-LHC Run 4 start}: June 2030.
\end{itemize}
\end{tcolorbox}

\section{Higgs Boson Mass: Record Precision}

ATLAS (Lepton-Photon Conference, Melbourne):
\begin{equation}
m_H = 125.11 \pm 0.11~\text{GeV} \quad (0.09\%~\text{precision})
\end{equation}
Combining diphoton ($125.22 \pm 0.14$ GeV) and four-lepton channels.

CMS:
\begin{equation}
m_H = 125.35 \pm 0.15~\text{GeV} \quad (0.1\%~\text{precision})
\end{equation}

\textbf{DFD compatibility}: DFD predicts $m_H$ via the quartic coupling from the spectral action. The predicted range is consistent with both measurements. \GREEN.

\section{BSM Searches: Still Null}

LHC Run 3 continues with no BSM signals. This is consistent with DFD, which predicts no new particles beyond the SM.

\section{Agent 9 Assessment: Grade B}

HL-LHC timeline crystallized with LS3 starting July 2026. Higgs mass precision improved. No BSM.

\section{New Literature (Agent 9)}

\begin{enumerate}
\item \textbf{ATLAS} (2025) --- $m_H = 125.11 \pm 0.11$ GeV, record precision. \NEW.
\item \textbf{CMS} (2025) --- $m_H = 125.35 \pm 0.15$ GeV.
\item \textbf{HL-LHC schedule update} (CERN, 2026) --- LS3 from July 2026. \NEW.
\item \textbf{LHC operation and HL-LHC} (arXiv:2505.03535) --- Comprehensive status.
\item \textbf{LHCb Run 3} (2025--2026) --- Lepton universality; anomalies dissolved.
\end{enumerate}


%=============================================================================
\part{Agent 10: Astrophysics}
\label{part:agent10}
%=============================================================================

\section{Mandate}

Wide binary update. MOND tests. Astrophysical tests of gravity modifications.

\section{\NEW: Quality Framework Rules Out MOND in Wide Binaries}

\begin{tcolorbox}[colback=blue!5!white, colframe=blue!60!black,
  title={MNRAS (2026): Quality Framework for Wide Binary Gravity Tests}]
A rigorous quality framework (MNRAS 547, 2026) demonstrates that:
\begin{enumerate}
\item Undetected close binaries, chance alignments, and projection effects can mimic MOND-like signals.
\item A systematic checklist of best practices is introduced.
\item When the framework is applied, \textbf{no evidence for MOND} is found.
\end{enumerate}

This is the strongest methodological argument yet against MOND in wide binaries.
\end{tcolorbox}

\section{Wide Binary Field: Converging Toward GR}

\begin{itemize}
\item \textbf{Quality framework} (MNRAS 547, 2026): No MOND evidence after proper filtering. \NEW.
\item \textbf{Cookson et al.\ (2026)}: GR preferred with realistic triple modelling.
\item \textbf{Orbital modelling sensitivity} (arXiv:2603.11015): Results depend on orbit modelling choices.
\item \textbf{Gaia DR4 expectations}: Astrometric time series with Keplerian orbital solutions expected only in DR4, which will be definitive.
\end{itemize}

\textbf{DFD position}: DFD predicts no MOND-like departure in wide binaries. The field is converging toward GR, consistent with DFD. Status remains \YELLOW{} pending DR4.

\section{Agent 10 Assessment: Grade B+}

Upgrade from B to B+ based on the quality framework paper converging the field toward GR.

\section{New Literature (Agent 10)}

\begin{enumerate}
\item \textbf{Quality framework} (MNRAS 547, 2026) --- No MOND evidence. \NEW.
\item \textbf{Cookson et al.\ (2026)}, arXiv:2504.07569 --- GR preferred.
\item \textbf{arXiv:2603.11015} (March 2026) --- Orbital modelling sensitivity.
\item \textbf{arXiv:2512.25002} (December 2025) --- Wide binary orbital momenta.
\item \textbf{Hernandez et al.\ (2024)} --- Critical review.
\end{enumerate}


%=============================================================================
\part{Agent 11: CMB}
\label{part:agent11}
%=============================================================================

\section{Mandate}

$A_L$ update with joint ACT+SPT+Planck. Simons Observatory expansion. CMB birefringence.

\section{$A_L$: Joint Measurement Tightens}

\begin{tcolorbox}[colback=yellow!5!white, colframe=yellow!60!black,
  title={Qu et al.\ (PRL 136, 2026): $A_L^{\text{recon}} = 1.025 \pm 0.017$}]

The joint ACT + SPT + Planck CMB lensing analysis achieves SNR $= 61$ (up from $\sim 40$ per experiment individually). The combined measurement:
\begin{equation}
A_L^{\text{recon}} = 1.025 \pm 0.017
\end{equation}

Also:
\begin{equation}
S_8^{\text{CMBL}} \equiv \sigma_8 (\Omega_m / 0.3)^{0.25} = 0.825^{+0.015}_{-0.013}
\end{equation}

\textbf{DFD prediction}: $A_L = 1.00$--$1.05$. The observed $1.025 \pm 0.017$ is \emph{exactly} in the DFD range. However, this prediction is ad hoc until Phase 0 provides a proper Boltzmann computation.

\textbf{Status}: \YELLOW{} (prediction exists but is not derived from Phase 0).
\end{tcolorbox}

\section{\NEW: Simons Observatory UK SATs Deployment}

\begin{itemize}
\item The SO Large Aperture Telescope (LAT) achieved first light on February 22, 2025 (Mars observation).
\item \textbf{SO:UK SATs}: Two new mid-frequency dichroic SATs (90 and 150 GHz) deploying in 2026. Different technology from existing SATs (UK-funded).
\item Full SO data acquisition and optimization underway for cross-correlation with LSS surveys.
\end{itemize}

\section{\NEW: Dark Energy and Lensing Anomaly Connection}

arXiv:2502.04641 (February 2026): Investigates the impact of the Planck lensing anomaly on dark energy constraints. Using varying $A_L$ values with DESI BAO and Pantheon+ supernovae, finds that the $A_L$ anomaly and the dynamical DE preference are correlated.

\textbf{DFD implication}: If the $A_L$ anomaly is resolved by DFD's distance-bias mechanism, this could simultaneously affect the dynamical DE significance. Phase 0 is needed to disentangle.

\section{CMB Birefringence}

Signal update (ScienceDaily, March 2026): Continued evidence for cosmic birefringence in CMB polarization data. DFD does not predict birefringence from the $\psi$-field (the field is scalar, not pseudo-scalar). If confirmed at $> 5\sigma$, this would require extension of DFD or identification of an independent source.

\textbf{Status}: Currently at $\sim 3\sigma$. Monitor.

\section{Agent 11 Assessment: Grade A-}

The joint $A_L$ measurement is exactly in DFD's predicted range. SO expansion on track. The birefringence signal is a potential concern but below threshold.

\section{New Literature (Agent 11)}

\begin{enumerate}
\item \textbf{Qu et al.\ (PRL 136, 2026)} --- Joint ACT+SPT+Planck $A_L = 1.025 \pm 0.017$. \NEW.
\item \textbf{SO:UK SAT deployment} (2026) --- Two new mid-frequency SATs. \NEW.
\item \textbf{arXiv:2502.04641} (February 2026) --- Dark energy and lensing anomaly. \NEW.
\item \textbf{CMB birefringence update} (March 2026) --- Signal continues at $\sim 3\sigma$.
\item \textbf{CMB-S4 forecasts} (arXiv:2505.22775) --- Non-thermal light relics.
\end{enumerate}


%=============================================================================
\part{Agent 12: Large-Scale Structure}
\label{part:agent12}
%=============================================================================

\section{Mandate}

Euclid pre-registration finalization. DESI DR2 published results. Dynamical DE status.

\section{Euclid Pre-Registration: FINALIZE by 20 June 2026}

\begin{tcolorbox}[colback=red!5!white, colframe=red!60!black,
  title={DEADLINE: Euclid Pre-Registration --- 20 June 2026}]

\textbf{Key dates}:
\begin{itemize}
\item \textbf{Euclid Q2 (Quick Release 2)}: 24 June 2026.
\item \textbf{Euclid DR1}: 21 October 2026 ($\sim$2100 deg$^2$).
\item \textbf{DFD pre-registration deadline}: 20 June 2026 (before Q2).
\end{itemize}

\textbf{What must be pre-registered}:
\begin{enumerate}
\item DFD distance-bias prediction for weak lensing shear power spectrum $C_\ell^{\kappa\kappa}$.
\item Predicted deviation from $\Lambda$CDM in the $(\Omega_m, \sigma_8)$ plane.
\item Null hypothesis: $\Lambda$CDM. Alternative: DFD distance-bias with specific $\Delta\psi(z)$ profile.
\item Statistical threshold for claiming detection: $3\sigma$ in blind analysis.
\end{enumerate}

\textbf{Status}: Draft exists. Requires Phase 0 output (SymBoltz) to populate numerical predictions. Without Phase 0, the pre-registration is qualitative only.

\textbf{Critical dependency}: SymBoltz Phase 0 must execute BEFORE the pre-registration can be finalized.
\end{tcolorbox}

\section{Euclid Preparation: Higher-Order Weak Lensing}

Euclid Collaboration (A\&A, March 2026): ``Toward a DR1 application of higher-order weak lensing statistics.'' The HOWLS (Higher-Order Weak Lensing Statistics) analysis pipeline is being developed for DR1, with Fisher forecasts for a Euclid-like setup.

\textbf{DFD relevance}: Higher-order statistics (peak counts, Minkowski functionals) may be more sensitive to DFD's distance-bias than two-point shear correlations alone.

\section{DESI DR2: Published Results}

\begin{tcolorbox}[colback=blue!5!white, colframe=blue!60!black,
  title={DESI DR2 BAO: Published in Phys.\ Rev.\ D (October 2025)}]

\textbf{Key numbers}:
\begin{itemize}
\item $> 14$ million galaxies and quasars, 3 years of data.
\item BAO precision: $\sim 0.24\%$ (factor $\sim 2$ improvement over DR1).
\item Lyman-$\alpha$ forest: 1.1\% (along) $\times$ 1.3\% (transverse) = 0.65\% isotropic.
\item Dynamical DE preference: $2.8$--$4.2\sigma$ (unchanged from DR1, does not diminish).
\item Mild $2.3\sigma$ tension between BAO and CMB in flat $\Lambda$CDM.
\end{itemize}

\textbf{DFD assessment}: The persistence of the dynamical DE signal at $2.8$--$4.2\sigma$ is important for DFD. DFD's distance-bias mechanism produces an \emph{effective} dynamical DE signature without introducing a scalar field. If DESI's signal is confirmed, DFD provides a natural geometric explanation.

\textbf{Caveat}: The signal is prior-dependent. Extending the $w_0, w_a$ prior bounds can shift the preference to $\Lambda$CDM.
\end{tcolorbox}

\section{\NEW: DESI DR2 BAO--CMB Tension}

The $2.3\sigma$ tension between DESI DR2 BAO distances and Planck CMB in flat $\Lambda$CDM is new. This may be:
\begin{enumerate}
\item A statistical fluctuation.
\item Evidence for dynamical DE.
\item Evidence for DFD's distance-bias (the bias modifies the BAO distance ladder differently from the CMB).
\end{enumerate}

Phase 0 is needed to compute the DFD prediction for this tension quantitatively.

\section{Agent 12 Assessment: Grade B+}

DESI DR2 results crystallized. Euclid timeline confirmed. The pre-registration deadline creates urgency for Phase 0.

\section{New Literature (Agent 12)}

\begin{enumerate}
\item \textbf{DESI DR2 BAO} (Phys.\ Rev.\ D, October 2025; arXiv:2503.14738) --- Published. \NEW.
\item \textbf{DESI DR2 Ly-$\alpha$} (arXiv:2503.14739) --- Lyman-alpha forest BAO. \NEW.
\item \textbf{Euclid HOWLS} (A\&A, March 2026) --- Higher-order WL statistics for DR1. \NEW.
\item \textbf{Giar\`e et al.\ (2025)}, EPJC --- Challenges DESI dynamical DE.
\item \textbf{Robust evidence paper} (2026) --- Joint DESI+ACT+SPT+Planck supports dynamical DE.
\item \textbf{DESI DR2 Guide} (March 2025) --- Official summary.
\end{enumerate}


%=============================================================================
\part{Agent 13: Dark Sector}
\label{part:agent13}
%=============================================================================

\section{Mandate}

$\Sigma m_\nu$ monitoring. JUNO timeline. Dark sector constraints.

\section{$\Sigma m_\nu$: Marginal but Alive}

As detailed in Agent 5 (Theory document):
\begin{itemize}
\item $\Sigma m_\nu < 0.064$ eV (DESI DR2 + Planck, $\Lambda$CDM).
\item $\Sigma m_\nu < 0.082$ eV (ACT DR6 + DESI DR2, $w_0 w_a$CDM).
\item DFD prediction: 0.061 eV.
\item Status: Marginal in $\Lambda$CDM, safe in $w_0 w_a$CDM.
\end{itemize}

\section{JUNO Timeline}

\begin{center}
\begin{tabular}{lll}
\toprule
\textbf{Year} & \textbf{Milestone} & \textbf{Status} \\
\midrule
2025 & Detector filling complete & DONE \\
2026 & Data taking begins & IN PROGRESS \\
2027--2028 & Mass ordering determination ($3\sigma$) & TARGET \\
2029--2030 & Precision $\Delta m^2$ measurements & TARGET \\
\bottomrule
\end{tabular}
\end{center}

JUNO's mass ordering determination will be critical for DFD. DFD predicts normal ordering. If JUNO finds inverted ordering, $\Sigma m_\nu$ shifts to $\geq 0.10$ eV, which is safely above cosmological bounds but changes the DFD neutrino mass predictions.

\section{\NEW: Coupled Dark Matter--Dark Energy}

arXiv:2503.10806 (March 2026): DESI results hint at coupled dark matter and dark energy. If confirmed, this could provide an alternative to DFD's distance-bias mechanism --- or it could be the phenomenological signature of DFD's $\psi$-field coupling to matter.

\section{Agent 13 Assessment: Grade B}

No change from i43. Neutrino mass situation marginal. JUNO data taking on track.

\section{New Literature (Agent 13)}

\begin{enumerate}
\item \textbf{DESI DR2 neutrino} (arXiv:2503.14744) --- $\Sigma m_\nu < 0.064$ eV. \NEW.
\item \textbf{arXiv:2503.10806} (March 2026) --- Coupled DM--DE from DESI. \NEW.
\item \textbf{arXiv:2603.10787} (March 2026) --- ACT DR6 + DESI DR2 neutrino. \NEW.
\item \textbf{JUNO status} (2026) --- Data taking commenced.
\item \textbf{Gaussian DE EoS} (MNRAS, 2026) --- Probing DE dynamics with DESI BAO.
\end{enumerate}


%=============================================================================
\part{Agent 14: SymBoltz / Numerics}
\label{part:agent14}
%=============================================================================

\section{Mandate}

\textbf{CRITICAL}: SymBoltz Phase 0 EXECUTION. This is the single highest priority of i44.

\section{SymBoltz.jl: Published and Ready}

\begin{tcolorbox}[colback=green!5!white, colframe=green!60!black,
  title={SymBoltz.jl Published in A\&A (March 2026)}]

\textbf{Reference}: Hersle (2026), A\&A; arXiv:2509.24740.

\textbf{Capabilities for DFD Phase 0}:
\begin{itemize}
\item Approximation-free Einstein--Boltzmann solver (no Limber, no tight-coupling switching).
\item Symbolic equation specification: DFD's modified Friedmann equations can be inserted directly.
\item Forward-mode automatic differentiation: Fisher forecasts for DFD parameters.
\item Replaceable physical components: DFD's $\psi$-field can be added as a module.
\item Flat Newtonian gauge, CDM, photons, baryons, RECFAST, massive neutrinos, $w_0 w_a$ DE.
\end{itemize}

\textbf{What remains}: Implementing DFD's specific modifications:
\begin{enumerate}
\item Distance-bias correction: $D_L^{\text{DFD}} = D_L^{\Lambda\text{CDM}} \cdot e^{\Delta\psi(z)}$.
\item Modified growth function: $f\sigma_8(z)$ with $\psi$-field contribution.
\item CMB power spectrum: $C_\ell^{TT}$, $C_\ell^{EE}$, $C_\ell^{\phi\phi}$ with distance-bias.
\end{enumerate}
\end{tcolorbox}

\section{Phase 0 Execution Plan}

\begin{tcolorbox}[colback=red!5!white, colframe=red!60!black,
  title={PHASE 0: 12-Hour Execution Plan}]

\textbf{Input}: SymBoltz.jl v1.0.0 + DFD $\psi$-field module (to be written).

\textbf{Steps}:
\begin{enumerate}
\item \textbf{Hour 0--2}: Install SymBoltz.jl. Verify $\Lambda$CDM baseline matches Planck 2018.
\item \textbf{Hour 2--5}: Implement DFD distance-bias as a custom component. Define $\Delta\psi(z)$ profile from DFD field equations.
\item \textbf{Hour 5--8}: Compute $C_\ell^{TT}$, $C_\ell^{EE}$, $C_\ell^{\phi\phi}$ for DFD. Compare to Planck.
\item \textbf{Hour 8--10}: Compute $A_L^{\text{DFD}}$ and $R_{31}^{\text{DFD}}$ (the two key observables).
\item \textbf{Hour 10--12}: Fisher forecast for Euclid DR1 sensitivity to DFD parameters. Generate pre-registration numbers.
\end{enumerate}

\textbf{Output}: Numerical predictions for $A_L$, $R_{31}$, $f\sigma_8(z)$, and Euclid forecast.

\textbf{Blocker}: None. SymBoltz.jl is published. Julia is available. The $\Delta\psi(z)$ profile is defined in DFD theory.

\textbf{Honest assessment}: This plan has been ``ready'' since i42. The programme must stop planning and start executing. Phase 0 is the single most important deliverable.
\end{tcolorbox}

\section{\NEW: gCAMB --- GPU-Accelerated Alternative}

arXiv:2509.25110; A\&C (2026): gCAMB is a GPU-accelerated version of CAMB achieving $\sim 100\times$ speedup. Potentially useful for MCMC sampling after Phase 0 provides the DFD modifications.

\section{Other Numerical Tools}

\begin{itemize}
\item \textbf{CAMB 1.6.5} (February 2026): Latest release. Stable.
\item \textbf{Lee et al.\ (2025)}: CLASS speedup techniques.
\end{itemize}

\section{Agent 14 Assessment: Grade A}

SymBoltz is published. Phase 0 plan is well-defined. The grade reflects readiness, not execution --- which has not yet occurred.

\textbf{Warning}: If Phase 0 is not executed by i45, Agent 14's grade will be downgraded to B regardless of readiness.

\section{New Literature (Agent 14)}

\begin{enumerate}
\item \textbf{SymBoltz.jl} (A\&A, March 2026; arXiv:2509.24740) --- Published. \NEW.
\item \textbf{gCAMB} (arXiv:2509.25110; A\&C, 2026) --- GPU-accelerated CAMB. \NEW.
\item \textbf{CAMB 1.6.5} (February 2026) --- Latest release.
\item \textbf{Lee et al.\ (2025)} --- CLASS speedup.
\item \textbf{CosmoCentral.jl} --- Julia cosmology infrastructure.
\end{enumerate}


%=============================================================================
\section*{End of i44 Observables Document}
%=============================================================================

\end{document}
